\chapter{Troubleshooting}
\label{ch:troubleshooting}

Start every investigation with the non-mutating facts:

\begin{lstlisting}[language=bash]
blackship capabilities
blackship check
blackship ps
blackship up --all --dry-run
blackship gc --dry-run
\end{lstlisting}

Add \command{--verbose} for command diagnostics. Preserve audit and scope data
before manual cleanup.

\section{A jail will not start}

\begin{enumerate}
  \item run \command{check} for dependency, name, and path errors;
  \item verify the selected release with \command{releases verify};
  \item inspect the dry-run path, address, hooks, and dataset action;
  \item confirm RCTL is enabled if resources are configured;
  \item inspect named network state and required kernel modules; and
  \item run \command{cleanup <jail>} only after recording the original error.
\end{enumerate}

A VNET or resource-limit error is fatal by design. Fix the host capability or
remove the unsupported requirement; do not expect Blackship to run the jail
with weaker isolation.

\section{RCTL is unavailable}

Check the read-only loader tunable:

\begin{lstlisting}[language=bash]
sysctl kern.racct.enable
\end{lstlisting}

It cannot be enabled at runtime. Add \command{kern.racct.enable="1"} to
\file{/boot/loader.conf} and reboot. After reboot, validate the actual rules
with \command{rctl -l jail:<full-name>}.

\section{Network creation or attachment fails}

\begin{lstlisting}[language=bash]
blackship network list
ifconfig -g blackship
sysctl net.inet.ip.forwarding
pfctl -sr
pfctl -sn
\end{lstlisting}

For netgraph, inspect node names and avoid reusing an \command{if\_bridge} name.
For epair, confirm the bridge and interface modules are available. Verify that
the jail address belongs to the named subnet and is not already leased.

If outbound packets leave the jail but not the host, check forwarding and NAT.
If inbound forwarding fails, confirm both the Blackship anchor and the anchor
load statements in \file{/etc/pf.conf}; then run \command{blackship setup} after
any PF reload.

\section{The jail has no DNS}

Inspect \file{/etc/resolv.conf} inside the jail and the manifest's DNS mode.
Custom nameservers must be reachable from the jail's subnet. In inherited mode,
ensure the host resolver file existed when the root was prepared. DNS failure
with working IP connectivity is not a bridge failure.

For an ephemeral \command{run} on a named network, Blackship removes host
loopback and interface-scoped link-local nameservers because its VNET stack
cannot reach them. If the host resolver contains only such entries, creation
fails with an instruction to supply a reachable server:

\begin{lstlisting}[language=bash]
blackship run --name probe --release 15.1-RELEASE \
  --network default --dns 192.0.2.53 -- host freebsd.org
\end{lstlisting}

Use repeated \command{--dns} options for redundancy. Do not point a VNET jail
at \command{127.0.0.1} unless it runs its own resolver there.

\section{A netgraph bridge name is already in use}

\command{ng\_ether} gives ethernet interfaces netgraph node names. If an
\command{if\_bridge} named \command{blackship0} already exists, create the
netgraph network as \command{ngbridge0} or another distinct name. The two bridge
types share a naming namespace even though their switching implementations
differ.

\section{EVA reports restart-required drift}

Path, VNET, addresses, and certain jail parameters are immutable for Blackship's
live-reconciliation model. EVA reports them without applying them. Schedule:

\begin{lstlisting}[language=bash]
blackship restart <jail> --dry-run
blackship restart <jail>
\end{lstlisting}

If \command{jls} cannot report a requested key, EVA may label it unreadable and
apply an updatable desired value blind. Review that result on the host's exact
FreeBSD version.

\section{A build unexpectedly reruns steps}

The layer key includes the base release, prior instructions, build arguments,
and copied content. A change early in the Jailfile correctly invalidates later
layers. Confirm that generated files in a \command{COPY} source are stable and
that the same build arguments are supplied. Check available ZFS snapshots; the
layer cache is a ZFS feature.

\section{Snapshot rollback is refused}

Stop the jail before rollback. If newer snapshots exist, a non-forced rollback
protects them. List snapshots, export anything needed, and only then decide
whether \command{--force} is appropriate.

\section{Residual resources remain after a crash}

Run a GC dry run and compare its report with \command{mount}, \command{jls},
\command{ifconfig -g blackship}, PF anchors, rctl rules, and ZFS datasets.
Normal GC reclaims proven ownership. Use targeted \command{cleanup} for one
jail. Force only after confirming that a reported resource does not belong to a
different live workload.

\section{Automatic escalation fails}

Run the command once with \command{--no-escalate} to confirm that privilege is
the only issue. Check \command{capabilities} for installed helper paths. For
mdo, verify the applicable \command{mac\_do} rule; installation alone does not
grant authorization. For unattended execution, ensure doas/sudo policy permits
the command without a password and use \command{--non-interactive}.

\section{JSON automation receives an error}

Confirm the global option precedes the command and that the command supports the
contract. Streaming health/watch and log/follow output is not a single JSON
envelope. On failure read stderr, check \setting{ok = false}, branch on the
stable error \setting{code}, and log the message for operators.

\section{Raw sockets or chflags inside an ephemeral run}

Ephemeral \command{run} jails are intentionally not granted raw sockets,
arbitrary socket families, or chflags. Use a managed manifest jail with an
explicitly reviewed policy for workloads that require those facilities. Do not
broaden all ephemeral jobs to accommodate one diagnostic command.

\section{When to report a bug}

Include the Blackship version, \command{capabilities} output, FreeBSD release,
minimal manifest with secrets removed, exact command, verbose error, relevant
audit records, and whether ZFS/PF/RCTL/VNET are enabled. State whether the issue
reproduces after a dry-run-only invocation. Never attach private keys, tokens,
full production roots, or unsanitized ZFS streams.
